\bookchapter{บทที่ 2: จัดระเบียบห้องดิจิทัลของคุณ — การจัดการไฟล์}

\chapterintro{โทรศัพท์ใช้นานๆ จะรก นิสัยที่ดีเริ่มจากการจัดการไฟล์ เหมือนห้องของคุณที่ต้องจัดเป็นระยะ พื้นที่ดิจิทัลของคุณก็ต้องดูแล โทรศัพท์ที่เป็นระเบียบไม่แค่หาของง่าย ยังทำงานลื่นไหลด้วย}

\booksection{ระบบจัดเก็บข้อมูลสามชั้น: คลังดิจิทัลของคุณ}

\begin{enumerate}
  \item \textbf{ตัวเครื่องจัดเก็บ}: โซนแกนกลาง ความเร็วเร็ว ระบบและแอปทั้งหมดอยู่ตรงนี้ เป็นพื้นที่ล้ำค่าที่สุดของโทรศัพท์ ต้องใช้อย่างช่างคำนวณ
  \item \textbf{SDการ์ด(ถ้ารองรับ)}: ตู้เสื้อส่วนขยาย ความเร็วช้ากว่านิดหน่อย เหมาะกับหนัง รูปภาพ และไฟล์ใหญ่ๆ ข้อสังเกต: ตอนนี้โทรศัพท์ฟลากชิปจำนวนมากไม่รองรับ SD การ์ดแล้ว
  \item \textbf{คลาวด์จัดเก็บ}: ตู้นิรภัย + เว็บไดรฟ์ซิงก์ กันของหาย และซิงก์ข้ามหลายอุปกรณ์ได้ iCloud, Baidu Netdisk, OneDrive ล้วนเป็นตัวเลือกยอดนิยม
\end{enumerate}

\booksection{รู้จัก “File Manager”}

แอนดรอยด์มักมีแอป [ไฟล์] หรือ [ไฟล์จัดการ] ส่วน iOS ใช้แอป [ไฟล์] ที่มาในตัว ทั้งสองคือ “ตัวจัดการทรัพยากร” ที่ไว้ดูไฟล์ทั้งหมดในโทรศัพท์

\booksubsection{ความหมายของโฟลเดอร์ที่พบบ่อย:}

\begin{itemize}
  \item \textbf{DCIM}: รูปและวิดีโอที่ถ่ายจากกล้อง
  \item \textbf{Pictures}: รูปอื่นๆ เช่นภาพหน้าจอ (สกรีนช็อต) และรูปที่บันทึกมา
  \item \textbf{Download}:เบราว์เซอร์ดาวน์โหลดของไฟล์
  \item \textbf{Music}:เพลงไฟล์
  \item \textbf{Movies}:วิดีโอไฟล์
  \item \textbf{Documents}:เอกสารไฟล์
  \item \textbf{Android}: โฟลเดอร์ระบบของแอนดรอยด์ ข้างในคือข้อมูลของแอปแต่ละตัว ปกติอย่าลบด้วยมือ
\end{itemize}

\booksection{โครงสร้างระบบไฟล์: เจาะลึก}

แอนดรอยด์อยู่บนฐาน Linux ระบบไฟล์จึงมีโครงสร้างลำดับชั้นที่แน่นอน แม้ผู้ใช้ทั่วไปไม่จำเป็นต้องเข้าใจลึก แต่รู้แนวคิดพื้นฐานบ้างก็มีประโยชน์

\booksubsection{พาร์ติชันตัวจัดเก็บ:}

\begin{itemize}
  \item \textbf{/system}: พาร์ติชันระบบ เก็บไฟล์ของระบบปฏิบัติการ ผู้ใช้ทั่วไปไม่มีสิทธิ์เข้าถึง ต้องใช้สิทธิ์ root
  \item \textbf{/data}: พาร์ติชันข้อมูล เก็บข้อมูลผู้ใช้และข้อมูลแอป เป็นพื้นที่ที่คุณใช้งานในชีวิตประจำวัน
  \item \textbf{/sdcard หรือ /storage/emulated/0}: พื้นที่จัดเก็บที่ผู้ใช้เข้าถึงได้ รูปของคุณ ไฟล์ที่ดาวน์โหลด ล้วนอยู่ตรงนี้
\end{itemize}

\booksubsection{Appข้อมูลจัดเก็บ:}

\begin{itemize}
  \item \textbf{ในเครื่องจัดเก็บ}: พื้นที่ส่วนตัวของแอป แอปอื่นเข้าไม่ได้ เมื่อถอนการติดตั้งแอปจะถูกลบไปด้วย
  \item \textbf{นอกเครื่องจัดเก็บ}: ไฟล์ที่แอปเก็บไว้ในพื้นที่ที่ผู้ใช้เข้าถึงได้ เช่นรูปภาพ ไฟล์ดาวน์โหลด เมื่อถอนการติดตั้งแอปอาจยังค้างอยู่ไม่ถูกลบ
\end{itemize}

\booksection{สำรองและกู้คืนข้อมูล: ทำประกันให้ข้อมูลของคุณ}

ข้อมูลไม่มีราคา, สำรองข้อมูลใช่ต้องเลี้ยงสำเร็จของคุ้นเคย

\booksubsection{สำรองข้อมูลกลยุทธ์:}

\begin{enumerate}
  \item \textbf{คลาวด์สำรองข้อมูล}: อัตโนมัติ สะดวก แต่ต้องใช้เครือข่ายและพื้นที่คลาวด์ iCloud, Huawei Cloud, Xiaomi Cloud ฯลฯ รองรับหมด
  \item \textbf{ภายในเครื่องสำรองข้อมูล}: ต่อคอมพิวเตอร์ด้วยสาย แล้วใช้เครื่องมือสำรองข้อมูลทางการ เร็ว พื้นที่เยอะ แต่ต้องทำด้วยมือเอง
  \item \textbf{คู่หนักสำรองข้อมูล}:คลาวด์สำรองข้อมูล+ภายในเครื่องสำรองข้อมูล, ที่สุดปลอดภัย
\end{enumerate}

\booksubsection{สำรองข้อมูลอะไร?}

\begin{itemize}
  \item \textbf{รูปภาพและวิดีโอ}: กินพื้นที่มากที่สุด และมีค่าที่สุดด้วย
  \item \textbf{สมุดโทรศัพท์}: เก็บไว้ในบัญชีคลาวด์สะดวกที่สุด
  \item \textbf{ประวัติแชท WeChat}: ใช้ฟังก์ชันสำรองข้อมูลของ WeChat เอง
  \item \textbf{สำคัญเอกสาร}: เก็บไว้บนคลาวด์ไดรฟ์หรือคอมพิวเตอร์
\end{itemize}

\booksection{ความปลอดภัยของไฟล์: ปกป้องความเป็นส่วนตัวของคุณ}

\begin{itemize}
  \item \textbf{ว่องไวรู้สึกไฟล์เข้ารหัส}: โทรศัพท์หลายรุ่นมีฟังก์ชัน "โฟลเดอร์ปลอดภัย" หรือ "พื้นที่ส่วนตัว" เอาไฟล์อ่อนไหวใส่ไว้ตรงนั้นได้ ต้องปลดล็อกด้วยรหัสผ่านหรือลายนิ้วก่อนถึงจะเข้าถึงได้
  \item \textbf{ลบไม่เท่ากับขายทำลาย}: การลบปกติเป็นแค่การทำเครื่องหมายว่าพื้นที่นี้เขียนทับได้ ข้อมูลยังอยู่ ถ้าจะลบให้สิ้นซากจริงๆ ต้องใช้เครื่องมือลบแบบปลอดภัยเฉพาะทาง หรือเขียนทับหลายรอบ
  \item \textbf{บนคลาวด์ปลอดภัย}: เลือกผู้ให้บริการคลาวด์ที่มีชื่อเสียง เปิดการยืนยันสองชั้น และตรวจสอบอุปกรณ์ที่ล็อกอินเป็นประจำ
\end{itemize}

\booksection{พื้นที่ไม่พอ? มาทำความสะอาดใหญ่}

\begin{itemize}
  \item \textbf{ล้างแคช}: ความจำชั่วคราวของแอป ล้างแล้วเป็นการเคลียร์ที่ ไปที่ ตั้งค่า > จัดเก็บ ล้างได้
  \item \textbf{ค้นหาใหญ่ไฟล์}: ใช้ตัวจัดการไฟล์จัดเรียงตามขนาด ตามล่า "ตัวร้าย" ที่กินพื้นที่ออกมา
  \item \textbf{ลบเหลือเหลือไว้}: หลังถอนการติดตั้งแอป บางทีจะมีไฟล์ขยะตกค้าง ล้างทิ้งด้วยมือ
  \item \textbf{ล้างอัลบั้ม}: ลบรูปที่ซ้ำกัน รูปเบลอ และภาพหน้าจอที่ไม่จำเป็น
  \item \textbf{ล้างวีแชต}: WeChat เป็นเจ้าที่กินพื้นที่ตัวยง การล้างประวัติแชทและแคชเป็นระยะจำเป็นมาก
\end{itemize}

\booksection{วิเคราะห์รูปแบบไฟล์ทั่วไปอย่างครบถ้วน}

ไฟล์ในโทรศัพท์มีนามสกุลมากมายจนเวียนหัว ทำความเข้าใจมันช่วยให้คุณจัดการไฟล์ได้ดีขึ้น:

\begin{center}
\begin{tabularx}{\linewidth}{|X|X|X|X|}
\hline
\textbf{รูปแบบ} & \textbf{ประเภท} & \textbf{คุณสมบัติ} & \textbf{แนะนำเปิดรูปแบบ} \tabularnewline
\hline
JPG/JPEG & รูปภาพ & บีบอัดแบบมีข้อมูลสูญหาย ไฟล์เล็ก ความเข้ากันได้ดีที่สุด เป็นรูปแบบเริ่มต้นของการถ่ายภาพบนโทรศัพท์ & แอปภาพ/ตัวดูรูปภาพใดก็ได้ \tabularnewline
PNG & รูปภาพ & บีบอัดแบบไม่สูญหาย รองรับพื้นหลังโปร่งใส ใช้บ่อยกับภาพหน้าจอและไอคอน ไฟล์ใหญ่กว่า JPG & แอปภาพ/ตัวดูรูปภาพใดก็ได้ \tabularnewline
MP3 & เสียง & บีบอัดแบบมีข้อมูลสูญหาย ไฟล์เล็ก ความเข้ากันได้ดีมาก เป็นรูปแบบเสียงที่ใช้ได้สากลที่สุด & โปรแกรมเล่นเพลงใดก็ได้ \tabularnewline
FLAC & เสียง & บีบอัดแบบไม่สูญหาย คุณภาพเสียงเทียบเท่า CD ไฟล์ใหญ่กว่า MP3 ถึง 3-5 เท่า ของโปรดสายออดิโอไฟแอน & โปรแกรมเล่นเพลง/แอปเล่นเฉพาะทาง \tabularnewline
WAV & เสียง & เสียงต้นฉบับไม่บีบอัด คุณภาพดีที่สุดแต่ไฟล์มโหฬาร ใช้บ่อยในงานอัดเสียงระดับมืออาชีพ & โปรแกรมเล่นเพลง \tabularnewline
M4A & เสียง & เข้ารหัสแบบ AAC คุณภาพเสียงดีกว่า MP3 ที่อัตราบิตเท่ากัน ใช้บ่อยในระบบนิเวศของแอปเปิล & Apple Music/โปรแกรมเล่นเพลง \tabularnewline
MP4 & วิดีโอ & รูปแบบครอบวิดีโอที่ใช้สากลที่สุด ความเข้ากันได้ดีมาก เป็นรูปแบบเริ่มต้นของวิดีโอที่อัดจากโทรศัพท์ & แอปภาพ/ตัวเล่นวิดีโอใดก็ได้ \tabularnewline
MKV & วิดีโอ & รูปแบบครอบแบบโอเพนซอร์ส รองรับหลายแทร็กเสียง หลายซับไตเติล ใช้บ่อยกับไฟล์หนัง & VLC/nPlayer \tabularnewline
DOC/DOCX & เอกสาร & เอกสาร Word DOCX เป็นรูปแบบบีบอัด ไฟล์เล็กกว่า DOC เป็นรูปแบบเวอร์ชันเก่า & WPS/Office \tabularnewline
PPT/PPTX & สาธิต & ไฟล์นำเสนอ PowerPoint PPTX เป็นรูปแบบบีบอัด & WPS/Office \tabularnewline
XLS/XLSX & ตาราง & สเปรดชีต Excel XLSX เป็นรูปแบบบีบอัด & WPS/Office \tabularnewline
PDF & เอกสาร & รูปแบบข้ามแพลตฟอร์ม จัดหน้าตายตัวแก้ไขไม่ได้ ใช้บ่อยกับสัญญา งานวิชาการ และหนังสืออิเล็กทรอนิกส์ & Adobe Reader/WPS \tabularnewline
ZIP & บีบอัด & รูปแบบบีบอัดที่ใช้สากลที่สุด ใช้ข้ามแพลตฟอร์มได้ดี & แอปในตัวระบบ/แอปแตกไฟล์ \tabularnewline
RAR & บีบอัด & อัตราการบีบอัดสูงกว่า ZIP แต่ต้องใช้ซอฟต์แวร์เฉพาะในการแตกไฟล์ & แอปแตกไฟล์ \tabularnewline
ISO & อิมเมจ & ไฟล์อิมเมจแผ่นดิสก์ คัดลอกเนื้อหาแผ่นมาทั้งหมด ใช้บ่อยกับแพ็กเกจติดตั้งระบบ & แอปแตกไฟล์/ไดรฟ์เสมือน \tabularnewline
HTML & หน้าเว็บ & ไฟล์หน้าเว็บ เปิดด้วยเบราว์เซอร์ หน้าเว็บที่บันทึกมักใช้รูปแบบนี้ & เบราว์เซอร์ \tabularnewline
\hline
\end{tabularx}
\end{center}

\begin{tcbpractice}{ฝึกหัด: จัดระเบียบโฟลเดอร์ "ดาวน์โหลด"}
เปิดตัวจัดการไฟล์ > เข้าไปที่โฟลเดอร์ “ดาวน์โหลด” ดูซิว่าข้างในมีอะไรบ้าง? อันที่ไม่ใช้ก็ลบทิ้ง อันที่ใช้ก็แยกหมวดหมู่ (เช่นสร้างโฟลเดอร์ “เอกสารงาน” “สติกเกอร์/มีม” ใส่เข้าไป) พลางสำรวจพื้นที่จัดเก็บดูว่าไฟล์ประเภทไหนกินที่มากที่สุด

\end{tcbpractice}

\booksection{การจัดการรูปภาพและระบบไฟล์: จัดสินทรัพย์ให้เป็นระเบียบ}

ตอนที่แล้วเราเรียน “การรู้จักโฟลเดอร์” กันไปแล้ว แต่แค่รู้จักไม่พอ ส่วนนี้จะพูดถึง\textbf{วิธีสร้างกฎการจัดการที่ทำได้จริงและใช้ได้ยาวนาน}ให้ไฟล์ของคุณหาเจอได้ทุกเมื่อ ไม่หาย ไม่ถูกแอบแอง

\booksubsection{การจัดหมวดหมู่ไฟล์และกฎการตั้งชื่อ}

รากของความโกลาหลคือการไม่มีกติกา การสร้างกติกาจริงๆ ง่ายมาก จุดสำคัญอยู่ที่\textbf{การทำต่ออย่างสม่ำเสมอ}.

\begin{itemize}
  \item \textbf{วิธีตั้งชื่อ "เวลา+เหตุการณ์"}: ชื่อโฟลเดอร์มีการประทับเวลาและคำอธิบายเหตุการณ์ ค้นกลับและนึกได้ง่ายในภายหลัง เช่น\texttt{202606\_ตั๋วการเดินทาง},\texttt{2025\_สรุปสิ้นปี},\texttt{2024\_วันเกิดลูก}แบบนี้เห็นทั้งไทม์ไลน์ และรู้เนื้อหาได้ในพริบตา
  \item \textbf{ระบบจำแนกสามระดับ}:
  \begin{itemize}
    \item \textbf{ระดับที่หนึ่ง}: แยกตามประเภทไฟล์ (เอกสาร / รูปภาพ / วิดีโอ / เสียง)
    \item \textbf{ระดับที่สอง}: แยกตามสถานการณ์ใช้งาน (งาน / ชีวิตประจำวัน / เรียนรู้)
    \item \textbf{ระดับที่สาม}: แยกตามเหตุการณ์หรือโครงการเฉพาะ (เช่น "2026Q2\_โครงการA")
  \end{itemize}
  \item \textbf{แผนล้างรายเดือน}: กำหนดวันคงที่หนึ่งวันทุกเดือน (แนะนำวันเงินเดือนออกหรือวันที่ 1 ของเดือน) จัดระเบียบโฟลเดอร์ "ดาวน์โหลด" รวมศูนย์ทีเดียว:\textbf{เก็บถาวรอันที่ใช้ได้ ลบอันที่ไม่ใช้ ย้ายเก็บอันที่ล้ำค่า}จัดเล็กๆ เดือนละครั้งตลอดปี สบายกว่ารวบทุกอย่างไว้กวาดใหญ่ทีเดียวมาก
\end{itemize}

\begin{tcbpractice}{รายการงานเดือนนี้}
เลือกสุดสัปดาห์ที่ว่างสักครั้ง แล้วทำตามลำดับนี้:

\begin{enumerate}
  \item สร้างในตัวจัดการไฟล์\texttt{2026Q2\_โครงการA}ในรูปแบบของโฟลเดอร์งานประจำเดือน
  \item ไล่ไฟล์สัปดาห์นี้ในโฟลเดอร์ "ดาวน์โหลด" ทั้งหมดหนึ่งรอบ เก็บที่ใช้ได้ จัดถาวร ลบที่ไม่ใช้
  \item สร้างให้รูปครอบครัวของปีนี้\texttt{2025ครอบครัว},\texttt{2026ครอบครัว}โฟลเดอร์รายปีแบบนี้
\end{enumerate}

\end{tcbpractice}

\booksubsection{กลยุทธ์จัดการไฟล์ซ้ำซ้อน}

โทรศัพท์ใช้นานจะ “อ้วนขึ้น” ส่วนใหญ่เป็นพวกไฟล์ซ้ำซ้อนที่แอบแองอยู่

\begin{itemize}
  \item \textbf{การล้างอัตโนมัติแบบ 72 ชั่วโมง}: App โซเชียลอย่าง WeChat QQ รับไฟล์ชั่วคราวมากมาย (สติกเกอร์/มีม รูปในแชท วิดีโอสั้น) ส่วนใหญ่ดูรอบเดียวก็ไม่ใช้แล้ว เข้าไปที่ [WeChat] > [ตั้งค่า] > [ทั่วไป] > [พื้นที่จัดเก็บ] เปิด "ล้างอัตโนมัติ" ตั้งแคชหมดอายุที่ 72 ชั่วโมง โทรศัพท์จะสะอาดขึ้นทันที แถมไม่ทิ้งบันทึกสำคัญ
  \item \textbf{การล้างเฉพาะจุดในพื้นที่จัดเก็บ}: แอนดรอยด์ [ตั้งค่า] > [จัดเก็บ] > [ล้างและเร่งความเร็ว] iOS [ตั้งค่า] > [ทั่วไป] > [พื้นที่จัดเก็บ iPhone] ระบบจะแยกหมวดให้ว่า "แคช / ไฟล์ชั่วคราว / แอปที่ถอนได้ / ไฟล์ใหญ่" จุดที่ควรล้าง: (1) ไฟล์แคช (2) แพ็กเกจติดตั้งหมดอายุ (แพ็กเกจรุ่นเก่าที่ค้างหลังอัปเดตแอป) (3) ไฟล์ซ้ำ (รูปชื่อและขนาดเหมือนกัน)
  \item \textbf{กฎการจัดหมวดหมู่อัตโนมัติ}: แอปอัลบั้มของแบรนด์ส่วนใหญ่รองรับจัดหมวดอัจฉริยะ: คน ทิวทัศน์ เอกสาร สกรีนช็อต เปิดไว้แล้ว รูปหมู่ที่ถ่ายใหม่จะถูกจัดเข้า "คน > ชื่อคน" สัญญาที่สแกนก็จะถูกจัดเข้า "เอกสาร" อัตโนมัติ\textbf{จับคู่กับการจัดระเบียบรายเดือน ช่วยลดโอกาส "หาของไม่เจอ" ได้อย่างมาก}.
\end{itemize}

\booksubsection{กฎการประมวลผลรูปภาพอย่างปลอดภัย}

รูปภาพเป็นประเภทข้อมูลที่ “อ่อนไหว” ที่สุดในโทรศัพท์ — ภาพถ่ายสบายๆ ภาพหนึ่งก็เปิดเผยตำแหน่ง อุปกรณ์ และเวลาของคุณได้ เราแนะนำให้\textbf{สร้างเวิร์กโฟลว์ “การป้องกันสองชั้น”}:

\begin{tcbwarn}{ขั้นตอนมาตรฐานในการส่งรูปบัตร/สกรีนช็อตที่อ่อนไหว}
\begin{enumerate}
  \item \textbf{ปิดบังข้อมูลที่มองเห็นได้ (มาร์กเอาต์)}: ใช้เครื่องมือในตัวระบบหรือภายนอก (เช่น Meitu, Snapseed)\textbf{ระบายทับเลขบัตรประชาชนบางหลักช่วงกลาง เลขทะเบียนรถ บ้านเลขที่ ชื่อ และจำนวนเงิน}และช่องข้อมูลสำคัญอื่นๆ
  \item \textbf{การลบข้อมูลเมตา}: ก่อนส่งให้กด "แชร์" หรือ "ส่งออก" ในอัลบั้ม ระบบจะถามว่าจะเก็บข้อมูลตำแหน่งไว้หรือไม่ ให้เลือก\textbf{"ไม่เก็บ"}แบบนี้รุ่นอุปกรณ์ที่ถ่าย พิกัด GPS และเวลาที่ถ่าย จะไม่ถูกส่งออกไปพร้อมรูป
  \item \textbf{ลายน้ำชั่วคราว}(ไม่บังคับ): ใส่ลายน้ำข้อความ "ใช้เพื่อดำเนินการ XX เท่านั้น" บนสกรีนช็อตที่อ่อนไหว เพื่อกันการถูกนำไปใช้ในทางที่ผิด
\end{enumerate}

\end{tcbwarn}

วิธีลบข้อมูลเมตา (EXIF) ที่ใช้กันทั่วไป (ขั้นพื้นฐาน ส่วนการป้องกันขั้นลึก ดูได้ที่บทที่ 6 หัวข้อ “การลบข้อมูลที่ระบุตัวตนออกจากรูปภาพ”):

\begin{itemize}
  \item \textbf{อยู่ในตัวระบบ}: ใน iOS [ตั้งค่า] > [ความเป็นส่วนตัวและความปลอดภัย] > [กล้อง] ปิด "ตำแหน่งที่แม่นยำ" ได้ ตอนแชร์เลือก "เฉพาะรูปที่เลือก" จะ\textbf{ลบ EXIF ส่วนใหญ่ออกโดยอัตโนมัติ}.
  \item \textbf{แอนดรอยด์}: แอปอัลบั้มของแบรนด์ส่วนใหญ่จะถามตอน "แชร์" ว่า "จะเอาข้อมูลตำแหน่งออกไหม" ติ๊กเลือกก็จบ
  \item \textbf{เครื่องมือระดับมืออาชีพ}: ExifTool (บนคอมพิวเตอร์) Metapho (iOS) Photo Metadata Remover (แอนดรอยด์) ลบข้อมูลเมตาได้หมดจริงๆ
\end{itemize}

สำหรับรูปที่ต้องการ\textbf{เก็บไว้ระยะยาว}ของรูปที่อ่อนไหว (บัตรประชาชน บัตรธนาคาร สำเนาทะเบียนบ้าน พาสปอร์ต) แนะนำให้เก็บเข้า\textbf{อัลบั้มเข้ารหัส}:

\begin{itemize}
  \item iPhone: [รูปภาพ] > เลือกรูป > ซ่อน (ดูได้หลังยืนยัน Face ID เท่านั้น)
  \item Huawei/Xiaomi: [อัลบั้ม] > กดค้างที่รูป > "ย้ายเข้าห้องนิรภัย"
  \item แอปเข้ารหัสจากภายนอก:\textbf{360 อัลบั้มปลอดภัย},\textbf{Tencent Album Manager (อัลบั้มส่วนตัว)},\textbf{KeepSafe}ฯลฯ มีรหัสผ่านป้องกันแยกของตัวเอง ซ่อนเร้นกว่า
\end{itemize}
